%Document Class
\documentclass[a4paper,12pt,reqno]{amsart}


  \headheight0.6in
  \headsep22pt
  \textheight23cm
  \topmargin-1.7cm
  \oddsidemargin 0.5cm
  \evensidemargin0.5cm
  \textwidth15.3cm

%Packages 
\usepackage{amscd,amssymb,amsmath,amsfonts,amsthm, mathrsfs,xspace}
\usepackage{graphicx}
\usepackage{verbatim,enumerate,paralist}
\usepackage{textcomp}
%\usepackage[authoryear]{natbib}
\usepackage[pagebackref,colorlinks]{hyperref}
\hypersetup{%
  urlcolor=blue,linkcolor=blue,citecolor=blue
}
\usepackage{pdfpages}
%\usepackage{pdfsync}%  remove this when finalised

% import these fonts by hand here to avoid clashes with usual blackboard bold usage.
%\pdfmapfile{+bbold.map}
\newcommand{\bbefamily}{\fontencoding{U}\fontfamily{bbold}\selectfont}
\newcommand{\textbbe}[1]{{\bbefamily #1}}
\DeclareMathAlphabet{\mathbbe}{U}{bbold}{m}{n}

% Macro to typeset part of a math formula at a bigger size
\def\makebigger#1#2#3{\hbox{#1$\mathsurround=0pt #2{#3}$}}
\def\bigger#1#2{{\relax\mathpalette{\makebigger#1}{#2}}}

% Macro for the category Delta
\newcommand{\DelCat}{{\bigger\large{\mathbbe{\Delta}}}}
\newcommand{\DelCAT}{{\bigger\Huge{\mathbbe{\Delta}}}}


\usepackage[T1]{fontenc}
\usepackage[all,2cell,poly,knot,tips]{xy}
\UseAllTwocells
% \CompileMatrices

\def\blue{\color{blue}}
\def\red{\color{red}}
\def\black{\color{black}}

%\usepackage{xypdf} % may not be necessary
\setlength{\marginparwidth}{2cm}
\newcommand{\missing}[1]{\marginpar{\color{red}missing:\\#1}}

%Theoretic Environment Setup
\newcounter{Definitioncount}
\setcounter{Definitioncount}{0}

\newtheorem*{Theorem}{Theorem}% no counter
\newtheorem{theorem}{Theorem}[section] % with counter
\newtheorem{lemma}[theorem]{Lemma}
\newtheorem{proposition}[theorem]{Proposition}
% \newtheorem*{Proposition}{Proposition}
\newtheorem{corollary}[theorem]{Corollary}
% \newtheorem*{Corollary}{Corollary}% no counter

\newtheorem*{Conjecture}{Conjecture}% no counter
\newtheorem*{Lemma}{Lemma}% no counter

\theoremstyle{definition}
\newtheorem*{Definition}{Definition}% no counter
\newtheorem*{Observation}{Observation}% no counter
\newtheorem*{Observations}{Observations}% no counter
\newtheorem*{Remark}{Remark}% no counter
\newtheorem{remark}[theorem]{Remark}
\newtheorem*{Example}{Example}% no counter
\newtheorem{example}[theorem]{Example}

\renewcommand{\theexample}{\arabic{example}}

\newtheoremstyle{fact}{\bigskipamount}{\medskipamount}{\upshape}{}{\itshape}{. }{ }{Fact}
\theoremstyle{fact}
\newtheorem*{Fact}{Fact}% no counter

\newtheoremstyle{genquest}{\bigskipamount}{\medskipamount}{\upshape}{}{\itshape}{. }{ }{General Question}
\theoremstyle{genquest}
\newtheorem*{GenQuest}{General Question}% no counter

% use \newtheoremstyle to underline the header
\newtheoremstyle{step}{2\bigskipamount}{\medskipamount}{\upshape}{}{\itshape}{. }{ }{\underline{Step~\thestep}}
\theoremstyle{step}
\newtheorem{step}{Step}[subsection]
\renewcommand{\thestep}{\arabic{step}}

\newcommand{\lra}{\longrightarrow}
\newcommand{\lla}{\longleftarrow}
\newcommand{\ra}{\rightarrow}
\newcommand{\la}{\leftarrow}
\newcommand{\Lra}{\Longrightarrow}
\newcommand{\Lla}{\Longleftarrow}
\newcommand{\ldual}[1]{\mathord{{\let\nolimits\relax\sideset{^\wedge}{}{#1}}}}
\newcommand{\laction}[2]{\mathord{{\let\nolimits\relax\sideset{^{#1}}{}{#2}}}}
\newcommand{\conj}[2]{\mathord{{\let\nolimits\relax\sideset{^{#1}}{}{#2}}}}
\newcommand{\xylongto}[2][2pt]{{\UseTips{}\xymatrix @C+#1{ \ar[r]^{#2}&}}}
\newcommand{\ola}{\overleftarrow}
\newcommand{\ora}{\overrightarrow}

\DeclareMathOperator{\im}{im}

%% Script Shortcuts 
\def\CA{{\mathscr A}}
\def\CB{{\mathscr B}}
\def\CC{{\mathscr C}}
\def\CD{{\mathscr D}}
\def\CE{{\mathscr E}}
\def\CF{{\mathscr F}}
\def\CG{{\mathscr G}}
\def\CH{{\mathscr H}}
\def\CI{{\mathscr I}}
\def\CJ{{\mathscr J}}
\def\CK{{\mathscr K}}
\def\CL{{\mathscr L}}
\def\CM{{\mathscr M}}
\def\CN{{\mathscr N}}
\def\CO{{\mathscr O}}
\def\CP{{\mathscr P}}
\def\CQ{{\mathscr Q}}
\def\CR{{\mathscr R}}
\def\CS{{\mathscr S}}
\def\CT{{\mathscr T}}
\def\CU{{\mathscr U}}
\def\CV{{\mathscr V}}
\def\CW{{\mathscr W}}
\def\CX{{\mathscr X}}
\def\CY{{\mathscr Y}}
\def\CZ{{\mathscr Z}}
\def\ordm{{\mathbf{m}}}
\def\ordn{{\mathbf{n}}}
\def\ordk{{\mathbf{k}}}


\newcommand{\ox}{\otimes}
\newcommand{\x}{\times}

\newcommand{\op}{\ensuremath{^{\textnormal{op}}}}


\def\theequation{{\thesection.\arabic{equation}}}


\begin{document}

\title{Combinatorial categorical equivalences}

\author{Stephen Lack}
\address{Department of Mathematics, Macquarie University, NSW 2109, Australia}
\email{steve.lack@mq.edu.au}

\author{Ross Street}
\address{Department of Mathematics, Macquarie University, NSW 2109, Australia}
\email{ross.street@mq.edu.au}

\thanks{Both authors gratefully acknowledge the support of the Australian Research Council Discovery Grant DP130101969; Lack acknowledges with equal gratitude the support of an Australian Research Council Future Fellowship}

\date{\today}
\maketitle

\tableofcontents

\section{Introduction}

The purpose is to generalize our work in \cite{PMSE} to include the 
Dold-Puppe-Kan Theorem \cite{Dold, DoldPuppe, Kan} 
on additive simplicial objects and chain complexes. 

\section{A setting for one direction}\label{settingone}

Let $\CP$ be a category. Let $\CZ$ be a class of non-invertible retractions
in $\CP$, closed under composition. Write $\bar{\CZ}$ for the left
ideal generated by $\CZ$. That is, $\bar{\CZ}$ consists of all the morphisms $g\circ z$ with $z\in \CZ$.  
Let $\CR$ be a class of morphisms $r$ in $\CP$ having the following properties: 
\begin{enumerate}[(a)]
\item $\CR \cap \bar{\CZ} = \emptyset$;
\item If $r : X \rightarrow A$ is in $\CR$ and $z:A\rightarrow U$ is in $\CZ$ then $z\circ r : X \rightarrow U$ is in $\bar{\CZ}$.
\end{enumerate}

\begin{example}\label{Ex1} As in \cite{PMSE}, take a category $\CA$ with a factorization system $(\CE,\CM)$. Assume that the pullback of any 
morphism with one in $\CM$ exist. Assume every morphism in $\CM$
is a monomorphism. Let $\CP$ be the category 
$\mathrm{Par}_{\CM}\CA$ of $\CM$-partial maps in $\CA$.
Take $\CZ$ to consist of the partial maps $m^*:A \lra U$ for
all non-invertible $m : U \lra A$ in $\CM$.
Let $\CR$ consist of the partial maps $e_* : X \lra Y$ 
for $e : X \lra Y$ in $\CE$.  

To see that property (a) holds, assume $e_* = g\circ m^*$ with 
$e\in \CE$ and $m\in \CM$. Then we have the diagram
\begin{equation}\label{pmcomp}
\begin{aligned}
\xymatrix{
& & A \ar[ld]^-{p_0} \ar@/_2pc/[lldd]_-{1_A} \ar[rd]_-{p_1} \ar@/^2pc/[rrdd]^-{e} \\
& U \ar[ld]^-{m} \ar[rd]_-{1_U} & & V \ar[rd]_-{g_1} \ar[ld]^-{g_0} \\
A & & U & & B}
\end{aligned}
\end{equation}
in which the diamond is a pullback. So $p_1$ is invertible and 
$mp_0=1_A$. Since $m$ is a monomorphism and a retraction, it is invertible.
So $m^*\notin \CZ$.

To see that the property (b) holds, consider $r = e_*$ and $z = m^*$
with $e\in \CE$ and $m\in \CM$. We have 
$$m^*\circ e_* = q_*\circ n^*$$ 
as per the pullback
\begin{equation}\label{ExAx2}
\begin{aligned}
\xymatrix{
P \ar[rr]^-{q} \ar[d]_-{n} && U \ar[d]^-{m} \\
X \ar[rr]_-{e} && A \ .}
\end{aligned}
\end{equation}
If $n$ is invertible then $mq=en\in \CE$, so $m\in \CE \cap \CM$
is invertible. So $m^*\in \CZ$ implies $n^*\in \CZ$.
\end{example}

\begin{example}\label{Ex2}
Take $\CP$ to be the category $\DelCat_{\bot,\top}$ of finite
non-empty ordinals $n = \{0,1,\dots,n-1\}$
with morphisms those functions which preserve
first element, last element and order.
Take $\CZ$ to consist of the non-invertible surjections $\sigma : m \lra n$
such that $\sigma (i) = 0$ implies $i = 0$; that is, $\sigma (1) = 1$.
In particular, for $k \neq 0$, the surjections $\sigma_k : m+1 \lra m$ with 
$\sigma_k(k)=\sigma_k(k+1)$ are in $\CZ$;
and every member of $\CZ$ is a composite of these.
Take $\CR$ to consist of the isomorphisms and the surjections 
$\sigma_0 : m+1 \lra m$ such that $\sigma_0 (1) = 0$.

To see that property (a) holds, suppose $\sigma_0 = \xi \sigma$ with
$\sigma : m+1 \lra n$ a non-invertible surjection. The
last statement implies $n< m+1$. Since $\sigma_0$ is surjective,
$\xi$ is surjective. So $n\geq m$. So $m=n$ and $\xi$ is an
identity. This means $\sigma = \sigma_0$ is not in $\CZ$.

To see that property (b) holds, just notice that, for $0<k$, 
$$\sigma_k \sigma_0 = \sigma_0 \sigma_{k+1} \ .$$ 
\end{example}

\bigskip
Returning to the general setting, we define a pointed category 
(that is, a category enriched in pointed sets) $\CD$ with the same objects
as $\CP$ and the non-zero morphisms those of $\CR$.
Composition is as in $\CP$ when the result is in $\CR$ but it is zero
otherwise.

Let $\CX$ be an additive category with sufficient limits including
finite direct sums. There is a functor
\begin{equation}\label{tilde}
\widetilde{(-)} : [\CP,\CX] \lra [\CD,\CX]_{\mathrm{pt}}
\end{equation}
where the codomain category consists of the pointed functors.
The definition on objects is:
\begin{equation}\label{tildeformula}
\widetilde{T}A = \bigcap_{(A\stackrel{z}\lra U) \in \CZ} {\mathrm{ker}}(Tz)
\end{equation}
and, for $r:A\lra B$ in $\CR$, the morphism 
$\widetilde{T}r : \widetilde{T}A \lra \widetilde{T}B$
is the restriction of $Tr : TA \lra TB$. 
Why does it restrict? Let $i_A : \widetilde{T}A \lra TA$
be the inclusion. Take any $z : B \lra V$ in $\CZ$.   
Then $(Tz)(Tr)i_A = T(zr)i_A = T(g\bar{z})i_A = (Tg)(T\bar{z})i_A = 0$
since $\bar{z} \in \CZ$ using property (b). 
So there exists $\widetilde{T}r$ such that $i_B (\widetilde{T}r) = (Tr)i_A$,
as claimed.

\section{A tighter setting}\label{settingtwo}

We again start with a category $\CP$.
We assume we have two classes $\CM$ and $\CR$ of morphisms in $\CP$. 

We ask that $\CM$ be closed under composition, contain the isomorphisms,
and consist of coretractions (= split monomorphisms). 

We ask that $\CR$ should contain the isomorphisms, contain composites
on both sides with isomorphisms, and have the property that 
$g\circ f \in \CR$ implies $g\in \CR$.

This time we {\em define} $\CZ$ to consist of those morphisms $z : X \rightarrow U$
which have a right inverse $z_* \in \CM$ and, if $g \circ z = f\circ m$ with 
$m\in \CM$, then $f \notin \CR$.
Clearly $\CZ$ includes no invertible morphisms yet is closed under composition.
In fact, if $g$ has a right inverse and $z\in \CZ$ then $g \circ z\in \CZ$. 

We require a set $\CI$ of pairs of morphisms $(m,m^*)$ where $m\in \CM$,
$m^*\circ m = 1$, and either $m^*\in\CZ$ or $m^*=1$. 
We assume every $z\in \CZ$ has the form $z = g\circ m^*$ for some 
retraction $g$ and $(m,m^*)\in \CI$ with $m^*\neq 1$.

The extra property we require is
\begin{align}\label{factorization}
& \textit{every } f\in \CP \textit{ factors uniquely as } f=m\circ r \circ z \textit{ with }
m\in \CM, r\in \CR, \nonumber \\ 
& \textit{ and either } z\in \CZ \textit{ or } z \textit{ an identity.} 
\end{align}

Property (a) of Section~\ref{settingone} is now a tautology.
We also have property (b) of Section~\ref{settingone}:
\begin{proposition}
Given $z : X \ra U$ in $\CZ$ and $r : A \ra X$ in $\CR$, there exist
$\zeta : A \ra B$ in $\CZ$ and some $g:B\ra U$ in $\CP$ such that 
$z\circ r = g\circ \zeta$. 
\end{proposition}
\begin{proof}
Apply \eqref{factorization} to $z\circ r$ to obtain $z\circ r = m\circ s \circ \zeta$
with $m\in \CM$, $s\in \CR$, and either $\zeta \in \CZ$ or $\zeta$ an identity.
If $\zeta$ is an identity then 
$m^*\circ z\circ r = s \in \CR$ where $m^*$ is a left inverse for $m$. 
It follows that $m^*\circ z \in \CR$. But this contradicts $z\in \CZ$. So
$\zeta$ is not an identity. Hence $\zeta \in \CZ$ and we take $g = m\circ s$.  
\end{proof} 

\begin{proposition}
Given $n : U \ra X$ in $\CM$ and $s : X \ra Y$ in $\CR$, there exist
unique $r : U \ra V$ in $\CR$ and $m:V\ra Y$ in $\CM$ such that 
$s\circ n = m\circ r$. 
\end{proposition}
\begin{proof}
Apply \eqref{factorization} to $s\circ n$ to obtain $s\circ n = m\circ r \circ z$
with $m\in \CM$, $r\in \CR$, and either $z \in \CZ$ or $z$ an identity.
By definition of $\CZ$ we cannot have both $z \in \CZ$ and $s\in \CR$.
So $z$ is an identity. So take $s\circ n = m\circ r$ as required.    
\end{proof}

It behooves us to show that our two examples of Section~\ref{settingone}
are still included. 

\begin{example}\label{Ex1+}
As in Example~\ref{Ex1}, take a category $\CA$ with a factorization 
system $(\CE,\CM)$. 
Assume that the pullback of any morphism with one in $\CM$ exist. 
Assume every morphism in $\CM$ is a monomorphism. 
Let $\CP$ be the category $\mathrm{Par}_{\CM}\CA$ of $\CM$-partial maps in $\CA$.
Identify $\CM$ in $\CA$ with the morphisms $m_*$ in $\CP$ for $m\in \CM$.
Take $\CR$ to consist of the morphisms $e_*$ in $\CP$ with $e\in \CE$.
With $\CZ$ as defined in this Section, for this example we claim
$$\CZ = \{m^* : m \text{ non-invertible in } \CM \} \ . $$
Take $z : X \lra U$ in $\CZ$. In order for a partial map $z = [z_0,V,z_1]$ to have 
a right inverse in $\CM$,
we see from the diagram
\begin{equation*}
\xymatrix{
& & U \ar[ld]^-{1_U} \ar@/_2pc/[lldd]_-{1_U} \ar[rd]_-{n} \ar@/^2pc/[rrdd]^-{1_U} \\
& U \ar[ld]^-{1_U} \ar[rd]_-{m} & & V \ar[rd]_-{z_1} \ar[ld]^-{z_0} \\
U & & X & & U}
\end{equation*}
with $m,n\in \CM$, we must have $n$, and hence $z_1$, invertible. So $z = m^*$.
So each $z\in \CZ$ is in the set claimed.

Now suppose we have $e_*\circ m_* = g \circ n^*$ with $e\in \CE$ and $m,n\in \CM$.
It follows that $n\circ g_0 = 1$. So $n^*\notin \CZ$.  
Our characterization of $\CZ$ follows.

We take $\CI = \{(m_*,m^*) : m\in \CM \}$.

Finally, every $f$ in $\CP$ uniquely factors as $f = m_*\circ e_* \circ n^*$ with
$m,n\in \CM$ and $e\in \CE$. But then $n^*$ is either invertible or in $\CZ$.    
\end{example}

\begin{example}\label{Ex2+}
As in Example~\ref{Ex2},take $\CP$ to be the category $\DelCat_{\bot,\top}$ of finite
non-empty ordinals $n = \{0,1,\dots,n-1\}$ with morphisms those functions which preserve
first element, last element and order.
We write $\sigma_k : m+1 \ra m$ for the order-preserving surjection which takes the value $k$ twice.
We write $\partial_k : m \ra m+1$ for the order-preserving injection which does not have $k$ in its image.
Take $\CM$ to consist of all the injective functions in $\DelCat_{\bot,\top}$. 
Take $\CR$ to consist of the isomorphisms and the surjections $\sigma_0 : m+1 \lra m$.
With $\CZ$ as defined in this Section, for this example we claim
$$\CZ = \{\tau \sigma_k : k > 0 \text{ and } \tau \text{ surjective} \} \ . $$
If $z$ is to have a right inverse it must be a surjection and not be $\sigma_0 : 2\ra 1$.
In fact no $\sigma_0 : m+1\ra m$ can be in $\CZ$ since $\partial_1 \sigma_0 = \sigma_0 \partial_2$
with $\sigma_0\in \CR$ and $\partial_2 \in \CM$. 

Now take any non-invertible surjection $\sigma$ with $\sigma(1) =1$.
We shall show that $\sigma \in \CZ$. Assume $\xi \sigma = \rho \partial$
where $\rho \in \CR$ and $\partial \in \CM$. Clearly $\rho$ cannot be
invertible or else $\sigma$ would be injective and so be invertible.
So $\rho = \sigma_0$. Using the fact that $\sigma \partial_i$ is an
identity for $i = 1$ and equal to $\partial_{i-1} \sigma_0$ for $i>1$,
we see that $\sigma_0 \partial$ is either injective or of the form
$\nu \sigma_0$ for some injective $\nu$. 
Factor $\xi = \mu \tau$ with $\tau$ surjective and $\mu$ injective. 
We have $\mu \tau \sigma$ either injective or equal to $\nu \sigma_0$.
The first case is impossible since $\sigma$ is not injective. By uniqueness
of factorization, the second case implies $\tau \sigma = \sigma_0$
and $\mu = \nu$. So $\tau = 1$ and $\sigma = \sigma_0$,
contradicting $\sigma(1) = 1$. So our assumption was impossible and
we are left with $\sigma \in \CZ$. In particular, this implies 
$\sigma_k \in \CZ$ for $k>0$ and hence $\tau \sigma_k \in \CZ$
for any surjective $\tau$.

Now suppose $(\sigma : m+1 \ra m) \in \CZ$. Let $0\le k < m$ be the largest 
such that $\sigma(k)=\sigma(k+1)$. Then $\sigma = \tau \sigma_k$ for some
surjection $\tau$. If $k=0$ then $\sigma = \sigma_0$ which we have seen is 
contrary to $\sigma \in \CZ$. So $\sigma = \tau \sigma_k$ with $k>0$.
This proves our characterization of $\CZ$.

We take $\CI = \{(\partial_k,\sigma_k) : k>0 \}$.
 
 Finally, for the factorization, take any $\xi$ in $\DelCat_{\bot,\top}$.
 If $\xi$ is injective, the required factorization is clear. 
 If $\xi = \partial \sigma_k$ with $\partial$ injective, the factorization is clear,
 breaking into two cases depending on whether $k=0$ or $k>0$.  
 Otherwise we can uniquely factorize $\xi = \partial \sigma_i \tau \sigma_k$ where
 $\partial$ is injective, $i<k$ and $\tau$ is a surjection for which $\tau(j)=\tau(j+1)$ implies $i<j<k$;
 again we have the desired factorization in each of the two cases $i=0$ or $i>0$.
 \end{example}

\bigskip
Returning to the general situation of this Section, notice that the formula \eqref{tildeformula}
for $\widetilde{T}: \CD\lra \CX$ now can be rewritten
\begin{equation}\label{tildeformula2}
\widetilde{T}A = \bigcap_{(U\stackrel{m}\lra A,m^*) \in \CI, m^*\neq 1} {\mathrm{ker}}(Tm^*) \ .
\end{equation}
We can also construct a functor
\begin{equation}\label{hat}
\widehat{(-)} : [\CD,\CX]_{\mathrm{pt}} \lra [\CP,\CX]
\end{equation}
whose value at the pointed functor $F : \CD\lra \CX$
is the functor $\widehat{F} : \CP \lra \CX$ defined as follows.   
On objects, put
$$\widehat{F}A = \sum_{(U\stackrel{m}\lra A,m^*) \in \CI}{FU} \ .$$
For a morphism $f:A\lra B$ in $\CP$, take the factorization
$f = n_1rz$ with $z:X\ra U$ in $\CZ$ or an identity, $r\in \CR$ 
and $n_1:V\ra Y$ in $\CM$.  
The morphism $\widehat{F}f: \widehat{F}A \lra \widehat{F}B$
has composite with the injection $\mathrm{in}_m : FU \ra \widehat{F}A$
equal to zero unless 
\begin{thebibliography}{00}

\bibitem{Berger1996} Clemens Berger, \textit{Op\'erades cellulaires et espaces de lacets it\'er\'es}, Annales de L'Institut Fourier \textbf{46(4)} (1996) 1125--1157.\label{Berger1996} 

\bibitem{CEF2012} Thomas Church, Jordan S. Ellenberg and Benson Farb, \textit{FI-modules: a new approach to stability for $S_n$-representations}, (\url{arXiv:1204.4533v2}, 28 Jun 2012).\label{CEF2012}

\bibitem{DaySt1995} Brian J. Day and Ross Street, \textit{Kan extensions along promonoidal functors}, Theory and Applications of Categories \textbf{1(4)} (1995) 72--77. \label{DaySt1995} 

\bibitem{DaySt1997} Brian J. Day and Ross Street, \textit{Monoidal bicategories and Hopf algebroids}, Advances in Math. \textbf{129} (1997) 99--157. \label{DaySt1997}

\bibitem{Dold} Albrecht Dold, \textit{Homology of symmetric products and other functors of complexes}, Annals of Math. \textbf{68} (1958) 54--80.\label{Dold}
 
\bibitem{DoldPuppe} Albrecht Dold and Dieter Puppe, \textit{Homologie nicht-additiver Funktoren. Anwendungen}, Ann. Inst. Fourier Grenoble \textbf{11} (1961) 201--312.\label{DoldPuppe}

\bibitem{FreydKelly} Peter J. Freyd and G. Max Kelly, \textit{Categories of continuous functors I}, J. Pure and Applied Algebra \textbf{2} (1972) 169--191; Erratum Ibid. \textbf{4} (1974) 121.\label{FreydKelly}

\bibitem{HillarDarren2007} Christopher J. Hillar and Darren L. Rhea, \textit{Automorphisms of finite abelian groups}, American Mathematical Monthly \textbf{114(10)} (2007) 917--923; \url{arXiv:math/0605185}.\label{HillarDarren2007}

\bibitem{Species} Andr\'e Joyal, \textit{Une theory combinatoire des series formelles}, Advances in Mathematics \textbf{42} (1981) 1--82.\label{Species}

\bibitem{AnalFunct} Andr\'e Joyal, \textit{Foncteurs analytiques et esp\`eces de structures}, Lecture Notes in Mathematics \textbf{1234} (Springer 1986) 126--159.\label{AnalFunct} 

\bibitem{TYBO} Andr\'e Joyal and Ross Street, \textit{Tortile Yang-Baxter operators in tensor categories}, J. Pure Appl. Algebra \textbf{71} (1991) 43--51.\label{TYBO}

\bibitem{ITD&QG} Andr\'e Joyal and Ross Street, \textit{An introduction to Tannaka duality and quantum groups}; Part II of Category Theory, Proceedings, Como 1990 (Editors A. Carboni, M.C. Pedicchio and G. Rosolini) Lecture Notes in Math. \textbf{1488} (Springer-Verlag Berlin, Heidelberg 1991) 411--492.\label{ITD&QG}

\bibitem{repfgl} Andr\'e Joyal and Ross Street, \textit{The category of representations of the general linear groups over a finite field}, J. Algebra \textbf{176 (3)} (1995) 908--946.\label{repfgl}

\bibitem{Kan} Daniel Kan, \textit{Functors involving c.s.s complexes}, Transactions of the American Mathematical Society \textbf{87} (1958) 330--346.\label{Kan}

\bibitem{PMSE} Stephen Lack and Ross Street, \textit{Morita equivalence of partial maps and strong epimorphisms}, preprint 2013.\label{PMSE}

\bibitem{Lindner1976} Harald Lindner, \textit{A remark on Mackey functors}, Manuscripta Mathematica \textbf{18} (1976) 273--278.\label{Lindner1976}

\bibitem{ML1963} Saunders Mac Lane, \textit{Natural associativity and commutativity}, Rice University Studies \textbf{49} (1963) 28--46.\label{ML1963}

\bibitem{NHA} Bodo Pareigis, \textit{A non-commutative non-cocommutative Hopf algebra in nature}, Journal of Algebra \textbf{70} (1981) 356--374.\label{NHA}

\bibitem{QG} Ross Street, \textit{Quantum Groups: A Path to Current Algebra}, Australian Mathematical Society Lecture Series \textbf{19} (Cambridge University Press, 2007).\label{QG} 

\bibitem{CBA} D. Tambara, \textit{The coendomorphism bialgebra of an algebra}, Journal of The Faculty of Science, The University of Tokyo, Section IA, Mathematics, \textbf{37} (1990) 425--456.\label{CBA}

\bibitem{Ulbrich} K.-H. Ulbrich, \textit{On Hopf algebras and rigid monoidal categories}, Israel Journal of 
Math. \textbf{72} (1990) 252--256.\label{Ulbrich}

\end{thebibliography}

\end{document}